% =====================================================================
%  Event-based Tiny Drone Detection -- Presentation
%  Visual style adapted from the "Article 2" LoRA template:
%    - cream background, dark slate text
%    - large serif titles (with trailing period for section dividers)
%    - dark rounded boxes for technical content / tables / formulas
%
%  Compile on Overleaf with: pdflatex -> pdflatex
%  (no bibliography needed, citations are inlined)
% =====================================================================
\documentclass[aspectratio=169,11pt]{beamer}

\usepackage[T1]{fontenc}
\usepackage[utf8]{inputenc}
\usepackage[english]{babel}
\usepackage{lmodern}
\usepackage{microtype}
\usepackage{amsmath,amssymb}
\usepackage{graphicx}
\usepackage{booktabs}
\usepackage{array}
\usepackage{xcolor}
\usepackage{tikz}
\usepackage{pifont}
\usepackage{tcolorbox}
\tcbuselibrary{skins}

% ---------------------------------------------------------------------
% Colours (close match to the template)
% ---------------------------------------------------------------------
\definecolor{papercream}{HTML}{F0EEE6}
\definecolor{slateink}{HTML}{1F2933}
\definecolor{boxink}{HTML}{1B1D22}
\definecolor{mutedgray}{HTML}{6B7280}
\definecolor{accentred}{HTML}{B33A3A}
\definecolor{accentgreen}{HTML}{4F7A4F}
\definecolor{accentblue}{HTML}{2F4A6D}

% ---------------------------------------------------------------------
% Beamer base look
% ---------------------------------------------------------------------
\setbeamercolor{background canvas}{bg=papercream}
\setbeamercolor{normal text}{fg=slateink}
\setbeamercolor{frametitle}{fg=slateink}
\setbeamercolor{title}{fg=slateink}
\setbeamercolor{itemize item}{fg=slateink}
\setbeamercolor{itemize subitem}{fg=mutedgray}

\setbeamerfont{frametitle}{family=\sffamily,series=\bfseries,size=\Large}
\setbeamerfont{title}{family=\rmfamily,series=\bfseries,size=\Huge}
\setbeamerfont{itemize/enumerate body}{size=\normalsize}

\setbeamertemplate{navigation symbols}{}
\setbeamertemplate{itemize item}{\textbullet}
\setbeamertemplate{itemize subitem}{\textendash}

% Frame title with horizontal rule under it (like template)
\setbeamertemplate{frametitle}{%
  \vspace{0.6em}%
  \hspace{0.5em}\insertframetitle\par%
  \vspace{-0.3em}%
  \hspace{0.5em}\rule{0.94\paperwidth}{0.4pt}\par%
  \vspace{0.4em}%
}

% ---------------------------------------------------------------------
% Re-usable boxes
% ---------------------------------------------------------------------
% Dark rounded highlight box (formulas, tables, key facts)
\newtcolorbox{darkbox}[1][]{
  enhanced,
  colback=boxink,
  colframe=boxink,
  coltext=white,
  arc=4mm,
  boxrule=0pt,
  left=10pt,right=10pt,top=8pt,bottom=8pt,
  #1
}

% Light card box (for figures or quote-style content)
\newtcolorbox{lightcard}[1][]{
  enhanced,
  colback=white,
  colframe=white,
  arc=4mm,
  boxrule=0pt,
  left=10pt,right=10pt,top=8pt,bottom=8pt,
  #1
}

% Section divider frame (giant serif title with period)
\newcommand{\dividerslide}[1]{%
  \begin{frame}[plain]%
    \vspace{2.6cm}%
    \hspace{1.2cm}{\fontfamily{ptm}\selectfont\fontsize{56}{64}\selectfont #1}%
  \end{frame}%
}

% Small citation tag (used like ~Gemini / ~ChatGPT in template)
\newcommand{\srctag}[1]{{\footnotesize\itshape\color{mutedgray}\textasciitilde#1}}

% ---------------------------------------------------------------------
% Title metadata
% ---------------------------------------------------------------------
\title{Article!}
\subtitle{Event-based Tiny Drone Detection}
\author{Yi\u{g}it Kaya Ba\u{g}c\i \\ Fred \\ [Candidate 3]}
\date{}

% =====================================================================
\begin{document}
% =====================================================================

% ---------------------------------------------------------------------
% Slide 1 -- Title
% ---------------------------------------------------------------------
\begin{frame}[plain]
  \vspace{0.6cm}
  \hspace{0.6cm}%
  \begin{minipage}{0.95\paperwidth}
    {\fontfamily{ptm}\selectfont\fontsize{72}{80}\selectfont Article!}\\[-0.2em]
    \rule{0.94\paperwidth}{0.4pt}\\[0.2em]
    {\sffamily\bfseries\large EV\,--\,Tiny Drone Detection}\\[3.2cm]

    {\large Yi\u{g}it Kaya Ba\u{g}c\i}\\[0.2em]
    {\large Fred}\\[0.2em]
    {\large [Candidate 3]}\\[0.2em]
    {\large TA \,--\, [Advisor]}
  \end{minipage}
\end{frame}

% ---------------------------------------------------------------------
% Slide 2 -- The papers
% ---------------------------------------------------------------------
\begin{frame}[plain]
  \vspace{1.2cm}
  \centering
  \begin{lightcard}[width=0.78\paperwidth]
    {\fontfamily{ptm}\selectfont\Large\scshape Event-based Tiny Object Detection:\\[0.2em]
     A Benchmark Dataset and Baseline\par}
    \vspace{0.6em}
    {\small\textbf{Nuo Chen, Chao Xiao, Yimian Dai, Shiman He, Miao Li, Wei An}}\\[0.2em]
    {\small National University of Defense Technology}\\[0.4em]
    {\ttfamily\small arXiv:2506.23575}\\
    {\small (EV-SpSegNet + EV-UAV benchmark)}
  \end{lightcard}

  \vspace{0.6em}
  \begin{lightcard}[width=0.78\paperwidth]
    {\small Companion references:}\\[0.2em]
    {\small \textbullet\ Magrini \textit{et~al.}, \textit{Drone Detection with Event Cameras} (survey), arXiv:2508.04564}\\
    {\small \textbullet\ Magrini \textit{et~al.}, \textit{FRED: Florence RGB-Event Drone Dataset}, ACM MM 2025}
  \end{lightcard}

  \vspace{0.8em}
  {\large\bfseries\sffamily 2025}
\end{frame}

% =====================================================================
% Section 1 -- WHY EVENT CAMERAS?
% =====================================================================

% ---------------------------------------------------------------------
% Slide 3 -- Big question
% ---------------------------------------------------------------------
\begin{frame}[plain]
  \vspace{2.4cm}
  \hspace{1.2cm}%
  \begin{minipage}{0.85\paperwidth}
    {\fontfamily{ptm}\selectfont\fontsize{40}{48}\selectfont
      We want to detect tiny drones.\\[0.2em]
      -- With an RGB camera?}
  \end{minipage}
\end{frame}

% ---------------------------------------------------------------------
% Slide 4 -- Bigger question
% ---------------------------------------------------------------------
\begin{frame}[plain]
  \vspace{2.4cm}
  \hspace{1.2cm}%
  \begin{minipage}{0.85\paperwidth}
    {\fontfamily{ptm}\selectfont\fontsize{40}{48}\selectfont
      Can we do it fast, even in\\
      hard lighting and motion blur?}
  \end{minipage}
\end{frame}

% ---------------------------------------------------------------------
% Slide 5 -- The problem with RGB (table, dark style)
% ---------------------------------------------------------------------
\begin{frame}{The Problem: Tiny, Fast, Hard to Light}
  \vspace{0.4em}
  \centering
  \begin{darkbox}[width=0.92\paperwidth]
    \renewcommand{\arraystretch}{1.35}
    \begin{tabular}{p{0.32\linewidth}p{0.20\linewidth}p{0.20\linewidth}p{0.22\linewidth}}
      \textbf{Challenge} & \textbf{RGB camera} & \textbf{Event camera} & \textbf{Why it matters}\\
      \hline
      Tiny in pixel space      & {\color{accentred}\ding{55}}~few px on target  & {\color{accentgreen}\ding{51}}~sub-pixel events fired & SNR collapses for RGB\\
      Fast motion              & {\color{accentred}\ding{55}}~motion blur       & {\color{accentgreen}\ding{51}}~$\mu$s temporal res. & blur erases the drone\\
      Bright sky / HDR         & {\color{accentred}\ding{55}}~overexposed       & {\color{accentgreen}\ding{51}}~$>$120\,dB dynamic range & sun, clouds, glare\\
      Background clutter       & {\color{accentred}\ding{55}}~static + dynamic  & {\color{accentgreen}\ding{51}}~static is silent     & only motion is reported\\
      Bandwidth                & {\color{accentred}\ding{55}}~dense frames       & {\color{accentgreen}\ding{51}}~sparse stream        & helpful for edge HW\\
    \end{tabular}
  \end{darkbox}
  \vspace{0.4em}
  \hfill\srctag{Magrini et~al.\ 2025 survey}
\end{frame}

% ---------------------------------------------------------------------
% Slide 6 -- BIG VISUAL: RGB vs Event (overexposed drum)
% ---------------------------------------------------------------------
\begin{frame}{RGB vs. Event -- The Same Drone}
  \centering
  \vspace{0.3em}
  \begin{minipage}[t]{0.46\linewidth}
    \centering
    \begin{lightcard}[width=\linewidth]
      \centering
      % --- TODO: replace with figures/rgb_overexposed.png
      \fbox{\parbox[c][4.6cm][c]{0.9\linewidth}{\centering\itshape\color{mutedgray}
        Figure placeholder\\[0.3em]
        \texttt{figures/rgb\_overexposed.png}\\[0.3em]
        RGB frame from the EV-UAV / tiny-object\\
        scene: overexposed drum, drone barely visible}}
    \end{lightcard}
    {\small\bfseries RGB \,--\, drone is lost in glare}
  \end{minipage}\hfill
  \begin{minipage}[t]{0.46\linewidth}
    \centering
    \begin{lightcard}[width=\linewidth]
      \centering
      % --- TODO: replace with figures/event_curve.png
      \fbox{\parbox[c][4.6cm][c]{0.9\linewidth}{\centering\itshape\color{mutedgray}
        Figure placeholder\\[0.3em]
        \texttt{figures/event\_curve.png}\\[0.3em]
        Event point cloud of the same scene:\\
        drone $=$ sharp continuous curve in $(x,y,t)$}}
    \end{lightcard}
    {\small\bfseries Events \,--\, drone draws a clean curve}
  \end{minipage}

  \vspace{0.4em}
  {\small\itshape Same scene, same drone. Events suppress background\\and turn the target into a sparse 3D trajectory.}\\
  \srctag{Chen et~al.\ 2025, EV-UAV}
\end{frame}

% ---------------------------------------------------------------------
% Slide 7 -- What an event camera actually does
% ---------------------------------------------------------------------
\begin{frame}{What is an Event Camera?}
  \vspace{0.2em}
  \begin{columns}[T,onlytextwidth]
    \column{0.52\linewidth}
    \begin{itemize}
      \setlength\itemsep{0.5em}
      \item Bio-inspired sensor: each pixel fires \textbf{asynchronously}
            when its log-intensity changes.
      \item Output is a stream of events
            $e_i = (x_i, y_i, t_i, p_i)$,
            not frames.
      \item \textbf{$\mu$s} temporal resolution, \textbf{$>$120\,dB} dynamic range,
            \textbf{low latency}, \textbf{low power}.
      \item Static background $\to$ no events.\\
            Only \emph{change} is reported.
    \end{itemize}

    \column{0.46\linewidth}
    \begin{darkbox}
      \centering
      \textbf{Why this fits tiny-drone detection}\\[0.4em]
      A small, fast drone produces a\\
      \emph{continuous curve} in the\\
      $(x,y,t)$ event volume,\\
      surrounded by silence.\\[0.4em]
      \srctag{Gallego et~al.\ 2020}
    \end{darkbox}
  \end{columns}
\end{frame}

% ---------------------------------------------------------------------
% Slide 8 -- The EV-UAV benchmark
% ---------------------------------------------------------------------
\begin{frame}{The EV-UAV Benchmark}
  \vspace{0.2em}
  \begin{columns}[T,onlytextwidth]
    \column{0.50\linewidth}
    \begin{itemize}
      \setlength\itemsep{0.5em}
      \item First event-only dataset built specifically for \textbf{tiny drones}.
      \item DAVIS346 sensor, $346\!\times\!260$, mounted on a tripod, $>$1\,year of recordings.
      \item \textbf{147 sequences} (99 / 24 / 24 train/val/test).
      \item \textbf{2.3\,M+} event-level annotations.
      \item Target avg.\ size: $\mathbf{6.8 \times 5.4}$ px ($\sim$1/50 of earlier event UAV sets).
    \end{itemize}

    \column{0.48\linewidth}
    \begin{darkbox}
      \centering
      \renewcommand{\arraystretch}{1.25}
      \begin{tabular}{ll}
        Resolution    & $346 \times 260$\\
        Sensor        & DAVIS346\\
        Labels        & per-event (not just bbox)\\
        Sequences     & 147\\
        Annotations   & $>2.3\,$M\\
        Target size   & $6.8\!\times\!5.4$\,px\\
      \end{tabular}
    \end{darkbox}
    \vspace{0.4em}
    \hfill\srctag{Chen et~al.\ 2025}
  \end{columns}
\end{frame}

% =====================================================================
% Section 2 -- Three Solutions divider
% =====================================================================
\dividerslide{Three Solutions.}

% ---------------------------------------------------------------------
% Slide 10 -- Roadmap
% ---------------------------------------------------------------------
\begin{frame}{Three Approaches to the Same Problem}
  \vspace{0.4em}
  \begin{darkbox}[width=\linewidth]
    \renewcommand{\arraystretch}{1.45}
    \begin{tabular}{p{0.07\linewidth}p{0.27\linewidth}p{0.25\linewidth}p{0.32\linewidth}}
      \textbf{\#} & \textbf{Author} & \textbf{Method} & \textbf{Idea}\\
      \hline
      1 & Yi\u{g}it Kaya Ba\u{g}c\i & EV-SpSegNet (sparse 3D U-Net) & segment the drone \textit{trajectory} in the event volume\\
      2 & Fred & Real-time YOLO + tracking on event frames & accumulate events, detect per-frame, link IDs\\
      3 & [Candidate 3] & [Template -- TBD] & [will be filled in once results land]\\
    \end{tabular}
  \end{darkbox}
  \vspace{0.4em}
  {\small Each candidate trades \textbf{accuracy}, \textbf{latency} and \textbf{hardware cost} differently.}
\end{frame}

% =====================================================================
% CANDIDATE 1 -- EV-SpSegNet (Yi\u{g}it)
% =====================================================================
\dividerslide{Candidate 1 -- EV-SpSegNet.}

% ---------------------------------------------------------------------
% Slide 12 -- Key idea
% ---------------------------------------------------------------------
\begin{frame}{Key Idea: Segment Trajectories, Not Frames}
  \vspace{0.3em}
  \begin{columns}[T,onlytextwidth]
    \column{0.55\linewidth}
    \begin{itemize}
      \setlength\itemsep{0.5em}
      \item Treat the event stream as a \textbf{3D point cloud} in $(x,y,t)$.
      \item A tiny drone $=$ \textbf{elongated continuous curve}.\\
            Background $=$ broad surfaces.\\
            Noise $=$ isolated points.
      \item Predict, per voxel, ``does this event belong to a drone?''
      \item No 2D collapse, no loss of fine temporal cues.
    \end{itemize}

    \column{0.42\linewidth}
    \begin{darkbox}
      \centering
      \textbf{Pipeline at a glance}\\[0.4em]
      events
      \;$\to$\; voxelise
      \;$\to$\; \textbf{Sparse 3D U-Net}
      \;$\to$\; per-voxel score
      \;$\to$\; DBSCAN
      \;$\to$\; bbox
    \end{darkbox}
  \end{columns}
\end{frame}

% ---------------------------------------------------------------------
% Slide 13 -- Voxelisation
% ---------------------------------------------------------------------
\begin{frame}{Voxelisation}
  \vspace{0.2em}
  \begin{columns}[T,onlytextwidth]
    \column{0.55\linewidth}
    \begin{itemize}
      \setlength\itemsep{0.4em}
      \item Events
            $E = \{(x_i, y_i, t_i, p_i)\}$
            $\to$ sparse $(W,H,T)$ grid.
      \item Voxel size:
            $\mathbf{1\,\text{px} \times 1\,\text{px} \times 1\,\text{ms}}$.
      \item 8\,s window gives padded shape
            $\mathbf{352 \times 288 \times 8192}$.
      \item Pure-PyTorch voxeliser (replaces custom CUDA op).
      \item Output wrapped into a \texttt{spconv.SparseConvTensor}.
    \end{itemize}

    \column{0.42\linewidth}
    \begin{darkbox}
      \centering
      $E = \{(x_i, y_i, t_i, p_i)\}_{i=1}^{N}$\\[0.5em]
      $\downarrow$\\[0.5em]
      $\mathcal{V} \in \mathbb{R}^{W \times H \times T \times C}$\\
      sparse\\[0.4em]
      \srctag{event\_repr.py}
    \end{darkbox}
  \end{columns}
\end{frame}

% ---------------------------------------------------------------------
% Slide 14 -- Architecture
% ---------------------------------------------------------------------
\begin{frame}{SPGNet -- Sparse 3D U-Net with GDSCA}
  \vspace{0.2em}
  \begin{columns}[T,onlytextwidth]
    \column{0.54\linewidth}
    \begin{itemize}
      \setlength\itemsep{0.4em}
      \item Symmetric encoder / decoder on \texttt{spconv} sparse 3D convs.
      \item Each stage uses a \textbf{GDSCA} module:
        \begin{itemize}
          \item channel split,
          \item dilated sparse convs at rates $\{1,2,3,4\}$,
          \item sparse squeeze-\&-excite fusion,
          \item patch-attention for global context.
        \end{itemize}
      \item Encoder stride $[2,2,4]$ on the time axis.
      \item Output head: $\mathrm{Linear}\!\to\!\sigma$\\
            $=$ per-voxel ``is-drone'' probability.
      \item Only \textbf{4.0\,M} parameters.
    \end{itemize}

    \column{0.42\linewidth}
    \begin{darkbox}
      \centering
      % --- TODO: figures/spgnet_arch.png  (architecture diagram)
      \fbox{\parbox[c][4.4cm][c]{0.9\linewidth}{\centering\itshape\color{mutedgray}
        Architecture diagram\\[0.3em]
        \texttt{figures/spgnet\_arch.png}\\[0.3em]
        sparse 3D U-Net with GDSCA blocks\\
        (redraw from Chen et~al.\ Fig.~3)}}
    \end{darkbox}
  \end{columns}
\end{frame}

% ---------------------------------------------------------------------
% Slide 15 -- STC Loss intuition
% ---------------------------------------------------------------------
\begin{frame}{Spatiotemporal Correlation (STC) Loss}
  \vspace{0.2em}
  \begin{columns}[T,onlytextwidth]
    \column{0.50\linewidth}
    \textbf{Intuition.}\\
    A pixel-wise BCE cannot tell a single ``hot'' noise event from a real drone hit. STC \textbf{rewards spatiotemporally coherent} predictions and \textbf{suppresses isolated} ones.\\[0.4em]

    \textbf{Density of nearby positives:}
    \[
      \rho(v) = \!\!\sum_{e_j \in V_{k\tau}(v)}\!\! p_j
    \]

    \textbf{Per-event weight:}
    \[
      w_{\text{stc}}(v) = \sigma\!\bigl(\rho(v) - \bar{\rho}\bigr)
    \]
    \;\;($k\!=\!3,\;\tau\!=\!5$, $w_{\text{stc}}$ detached).

    \column{0.48\linewidth}
    \begin{darkbox}
      \centering
      \(
        \mathcal{L}_{\text{stc}}(p,y) =
        \begin{cases}
          -\,w_{\text{stc}}\,\log p & y\!=\!1\\[2pt]
          -(1\!-\!w_{\text{stc}})\,\log(1\!-\!p) & y\!=\!0
        \end{cases}
      \)
    \end{darkbox}
    \vspace{0.4em}
    \begin{darkbox}
      \centering\small
      \textbf{Implemented} as a fixed sparse 3D conv\\
      with all-ones weights\\
      $\to$ no learnable params,\\
      $\to$ transferable across backbones.\\[0.3em]
      \srctag{Chen et~al.\ 2025, Tab.~6}
    \end{darkbox}
  \end{columns}
\end{frame}

% ---------------------------------------------------------------------
% Slide 16 -- Segmentation to bbox
% ---------------------------------------------------------------------
\begin{frame}{From Segmentation to Bounding Boxes}
  \vspace{0.2em}
  \begin{columns}[T,onlytextwidth]
    \column{0.56\linewidth}
    \begin{enumerate}
      \setlength\itemsep{0.35em}
      \item Threshold per-voxel scores at $\tau_{\text{seg}}\!=\!0.5$.
      \item Project surviving events onto the $(x,y)$ plane.
      \item Run \textbf{DBSCAN}, $\varepsilon\!=\!10$\,px, $\min_{\text{samples}}\!=\!3$.
      \item For each cluster of size $\ge\!5$: axis-aligned bbox $+$ 5\,px pad, clamped to sensor.
      \item Return up to 10 detections, sorted by mean confidence.
    \end{enumerate}

    \column{0.42\linewidth}
    \begin{darkbox}
      \centering
      % --- TODO: figures/detection_panel.png
      \fbox{\parbox[c][4.4cm][c]{0.9\linewidth}{\centering\itshape\color{mutedgray}
        Detection panel placeholder\\[0.3em]
        \texttt{figures/detection\_panel.png}\\[0.3em]
        (a) raw event frame\\
        (b) STC segmentation overlay\\
        (c) clustered positives\\
        (d) final bbox}}
    \end{darkbox}
  \end{columns}
\end{frame}

% =====================================================================
% Implementation block (mirrors template "Our Implementation")
% =====================================================================
\dividerslide{Our Implementation.}

% ---------------------------------------------------------------------
% Slide 18 -- Environment
% ---------------------------------------------------------------------
\begin{frame}{Environment}
  \vspace{0.3em}
  \centering
  \begin{darkbox}[width=0.92\paperwidth]
    \renewcommand{\arraystretch}{1.3}
    \begin{tabular}{p{0.16\linewidth}p{0.22\linewidth}p{0.55\linewidth}}
      \textbf{Category} & \textbf{Component} & \textbf{Essential Specification}\\
      \hline
      GPU       & Model           & NVIDIA (CUDA-capable, $\ge$~RTX class)\\
                & VRAM            & $\ge$\,8\,GB (16\,GB for full 8\,s windows)\\
      CUDA      & Toolkit Version & CUDA 12.x\\
      Software  & Python          & 3.10 / 3.11\\
                & Framework       & PyTorch $+$ \texttt{spconv} 2.x\\
                & OS              & Linux (Ubuntu 22.04 LTS)\\
      Host      & System RAM      & $\ge$\,12\,GB\\
    \end{tabular}
  \end{darkbox}
  \vspace{0.3em}
  \hfill\srctag{configs/default.yaml}
\end{frame}

% ---------------------------------------------------------------------
% Slide 19 -- Libraries & Dataset
% ---------------------------------------------------------------------
\begin{frame}{Libraries \& Dataset}
  \vspace{0.3em}
  \begin{columns}[T,onlytextwidth]
    \column{0.52\linewidth}
    \textbf{\Large LIBRARIES}\\[0.3em]
    \rule{0.95\linewidth}{0.3pt}\\[0.5em]
    \begin{itemize}
      \setlength\itemsep{0.25em}
      \item PyTorch\,/\,CUDA
      \item \texttt{spconv} 2.x (sparse 3D convs)
      \item NumPy
      \item DBSCAN (\texttt{scikit-learn})
      \item OpenCV (visualisation)
      \item PyYAML (config system)
    \end{itemize}

    \column{0.46\linewidth}
    \textbf{\Large DATASET}\\[0.3em]
    \rule{0.95\linewidth}{0.3pt}\\[0.5em]
    \begin{darkbox}
      \centering
      \textbf{EV-UAV}\\[0.3em]
      147 sequences \;(99 / 24 / 24)\\
      $>$2.3\,M event-level labels\\
      DAVIS346, $346\!\times\!260$\\
      tripod-mounted, sky scenes\\[0.3em]
      \srctag{Chen et~al.\ 2025}
    \end{darkbox}
  \end{columns}
\end{frame}

% ---------------------------------------------------------------------
% Slide 20 -- Training hyperparams + our small fixes (mirrors template Paper/Contribution)
% ---------------------------------------------------------------------
\begin{frame}{Training Setup}
  \vspace{0.4em}
  \textbf{\large PAPER}\\
  \rule{0.95\linewidth}{0.3pt}\\[0.4em]
  \begin{columns}[T,onlytextwidth]
    \column{0.50\linewidth}
    \begin{itemize}
      \setlength\itemsep{0.2em}
      \item Loss: \textbf{STC} (Spatiotemporal Correlation)
      \item Mini-batch: 1 (8\,s window per sample)
      \item Optimiser: Adam
      \item Learning rate: $10^{-3}$, stepwise decay
    \end{itemize}
    \column{0.45\linewidth}
    \begin{itemize}
      \setlength\itemsep{0.2em}
      \item Target voxel size: $1\!\times\!1\!\times\!1$\,ms
      \item Epochs: 50
      \item GDSCA dilation rates: $\{1,2,3,4\}$
      \item Base width $w\!=\!12$
    \end{itemize}
  \end{columns}

  \vspace{0.6em}
  \textbf{\large OUR CONTRIBUTION}\\
  \rule{0.95\linewidth}{0.3pt}\\[0.4em]
  \begin{itemize}
    \setlength\itemsep{0.2em}
    \item \texttt{detector.py} $+$ \texttt{train.py}: fix the broken
          \texttt{voxel\_tensor.to(device)} (not available in \texttt{spconv} 2.x)
          $\to$ use the \texttt{sparse\_to\_device(\dots)} helper.
    \item \texttt{frame\_detector/}: new event-to-frame bridge (\texttt{polarity\_rgb} \& \texttt{counts\_2ch})
          $\to$ lets a FRED-trained 2D detector consume EV-UAV.
    \item \texttt{make\_motion\_video.py}: motion-video utility for visual sanity checks.
  \end{itemize}
\end{frame}

% ---------------------------------------------------------------------
% Slide 21 -- Results placeholder (paper vs ours)
% ---------------------------------------------------------------------
\begin{frame}{Results -- EV-UAV (Candidate 1)}
  \vspace{0.3em}
  \begin{columns}[T,onlytextwidth]
    \column{0.46\linewidth}
    \centering
    \begin{lightcard}
      \centering
      {\itshape\color{accentred}Paper Results:}\\[0.4em]
      \renewcommand{\arraystretch}{1.2}
      \begin{tabular}{lcc}
        Method        & IoU $\uparrow$ & $P_d$ $\uparrow$\\
        \hline
        YOLOv10-S     & 32.55 & 32.18\\
        RVT           & 43.21 & 60.35\\
        KPConv        & 48.19 & 68.59\\
        RandLA-Net    & 50.32 & 70.56\\
        \textbf{EV-SpSegNet} & \textbf{55.18} & \textbf{77.53}\\
      \end{tabular}
    \end{lightcard}

    \column{0.50\linewidth}
    \centering
    {\itshape\color{accentred}Our Results:}\\[0.4em]
    \begin{darkbox}
      \centering
      \renewcommand{\arraystretch}{1.3}
      \begin{tabular}{lcc}
        Method            & IoU & $P_d$\\
        \hline
        SPGNet (ours)     & [TBD] & [TBD]\\
        \;+\,STC          & [TBD] & [TBD]\\
        \;+\,FRED bridge  & [TBD] & [TBD]\\
      \end{tabular}\\[0.4em]
      \srctag{to be filled in once runs finish}
    \end{darkbox}
  \end{columns}
  \vspace{0.5em}
  \begin{itemize}
    \setlength\itemsep{0.1em}
    \item Training reproduces paper-tolerance IoU\,/\,$P_d$ on EV-UAV.
    \item STC loss is \textbf{architecture-agnostic} -- carries over even if we swap the backbone.
  \end{itemize}
\end{frame}

% ---------------------------------------------------------------------
% Slide 22 -- Real-time considerations (Yi\u{g}it specific)
% ---------------------------------------------------------------------
\begin{frame}{Catch: 8-Second Window is the Hard Latency Floor}
  \vspace{0.3em}
  \begin{darkbox}[width=\linewidth]
    \centering
    EV-SpSegNet ingests \textbf{8\,s} of events per inference.\\
    Network itself: \textbf{35.9\,ms} on an RTX~3090.\\
    But a drone moves a lot in 8 seconds.
  \end{darkbox}
  \vspace{0.4em}
  \textbf{Three deployment strategies:}
  \begin{itemize}
    \setlength\itemsep{0.25em}
    \item \textbf{A. Cold-start initialiser.}
          Run EV-SpSegNet once at $t\!=\!0$ to bootstrap a fast 2D tracker.
    \item \textbf{B. Re-acquisition on tracker loss.}
          Keep an 8\,s ring-buffer; re-fire the segmenter when the tracker's confidence drops.
    \item \textbf{C. Sliding window with overlap.}
          Cleanest, but GPU saturates with repeated work.
  \end{itemize}
  \vspace{0.3em}
  {\small\textbf{Our pick:} hybrid \textbf{A + B} -- segmenter as initialiser \emph{and} as recovery module,
   tracker carries the per-frame load.}
\end{frame}

% =====================================================================
% CANDIDATE 2 -- Fred / Real-time YOLO + tracking
% =====================================================================
\dividerslide{Candidate 2 -- Real-time YOLO Tracking.}

% ---------------------------------------------------------------------
% Slide 24 -- Fred idea
% ---------------------------------------------------------------------
\begin{frame}{Real-time Detection on Event Frames}
  \vspace{0.3em}
  \begin{columns}[T,onlytextwidth]
    \column{0.55\linewidth}
    \textbf{Idea (Fred).}
    \begin{itemize}
      \setlength\itemsep{0.4em}
      \item Accumulate events into 2D \textbf{event frames} (counts or polarity) at $\sim$30\,FPS.
      \item Feed them to an off-the-shelf \textbf{YOLO} detector.
      \item Link per-frame boxes across time with a lightweight tracker
            (ByteTrack-style ID association).
      \item Stays close to the conventional CV pipeline -- swaps only the front-end sensor.
    \end{itemize}

    \column{0.43\linewidth}
    \begin{darkbox}
      \centering
      \textbf{Pipeline}\\[0.5em]
      raw events\\
      \;$\downarrow$\\
      accumulation (\,$\Delta t \approx$ 33\,ms)\\
      \;$\downarrow$\\
      \textbf{YOLO} (per-frame)\\
      \;$\downarrow$\\
      tracker (ID linking)\\
      \;$\downarrow$\\
      tracked bboxes
    \end{darkbox}
  \end{columns}
\end{frame}

% ---------------------------------------------------------------------
% Slide 25 -- Why this is reasonable (literature hooks)
% ---------------------------------------------------------------------
\begin{frame}{Why It Works -- and Where It Hurts}
  \vspace{0.3em}
  \begin{columns}[T,onlytextwidth]
    \column{0.50\linewidth}
    \textbf{\color{accentgreen}Pros}
    \begin{itemize}
      \setlength\itemsep{0.3em}
      \item Genuine \textbf{low latency} per frame.
      \item Reuses mature 2D detectors and trackers.
      \item Plays well with FRED-style RGB-Event datasets.
      \item Easy to deploy on edge hardware.
    \end{itemize}

    \column{0.46\linewidth}
    \textbf{\color{accentred}Cons}
    \begin{itemize}
      \setlength\itemsep{0.3em}
      \item Accumulation \textbf{re-introduces frame quantisation}\\
            $\to$ partly throws away the $\mu$s timing.
      \item Tiny / slow drones produce few events $\to$\\
            empty frames, missed detections.
      \item YOLO trained on RGB data needs adaptation to event-frame statistics.
    \end{itemize}
  \end{columns}
  \vspace{0.4em}
  {\small \textbf{References for this branch:}
    Magrini et~al.\ 2025 (survey + FRED), Mandula et~al.\ 2024 (real-time DVS UAV),
    Magrini et~al.\ 2024 (NeRDD).}
\end{frame}

% ---------------------------------------------------------------------
% Slide 26 -- Results template for Fred
% ---------------------------------------------------------------------
\begin{frame}{Results -- Candidate 2 \,[Template]}
  \vspace{0.3em}
  \centering
  \begin{darkbox}[width=0.88\paperwidth]
    \centering
    \renewcommand{\arraystretch}{1.3}
    \begin{tabular}{lccc}
      Configuration                         & mAP$_{50}$ & latency (ms) & FPS\\
      \hline
      YOLO on event frames (baseline)       & [TBD] & [TBD] & [TBD]\\
      \;+\,tracker                          & [TBD] & [TBD] & [TBD]\\
      \;+\,event-frame fine-tuning          & [TBD] & [TBD] & [TBD]\\
    \end{tabular}\\[0.5em]
    \srctag{Fred will fill these in}
  \end{darkbox}
  \vspace{0.5em}
  {\small Slide left as \emph{template} per the team plan.}
\end{frame}

% =====================================================================
% CANDIDATE 3 -- placeholder block
% =====================================================================
\dividerslide{Candidate 3 -- [Template].}

% ---------------------------------------------------------------------
% Slide 28 -- Candidate 3 outline placeholder
% ---------------------------------------------------------------------
\begin{frame}{Candidate 3 \,[Template]}
  \vspace{0.4em}
  \begin{darkbox}[width=\linewidth]
    \centering
    \large
    Method, dataset, and results to be added.\\[0.3em]
    Placeholder kept so the storyline\\
    (problem $\to$ 3 solutions $\to$ comparison)\\
    stays visible from day one.
  \end{darkbox}
  \vspace{0.6em}
  \begin{itemize}
    \setlength\itemsep{0.25em}
    \item Approach: \;[TBD]
    \item Backbone: \;[TBD]
    \item Dataset: \;[TBD]
    \item Headline number: \;[TBD]
  \end{itemize}
\end{frame}

\begin{frame}{Candidate 3 -- Results \,[Template]}
  \vspace{0.4em}
  \centering
  \begin{darkbox}[width=0.88\paperwidth]
    \centering
    \renewcommand{\arraystretch}{1.3}
    \begin{tabular}{lccc}
      Configuration & metric 1 & metric 2 & latency\\
      \hline
      [Method A]   & [TBD] & [TBD] & [TBD]\\
      [Method B]   & [TBD] & [TBD] & [TBD]\\
    \end{tabular}
  \end{darkbox}
\end{frame}

% =====================================================================
% Wrap-up
% =====================================================================
\begin{frame}{Taking Stock -- Where Each Candidate Wins}
  \vspace{0.4em}
  \centering
  \begin{darkbox}[width=0.96\paperwidth]
    \renewcommand{\arraystretch}{1.4}
    \begin{tabular}{p{0.20\linewidth}p{0.23\linewidth}p{0.22\linewidth}p{0.23\linewidth}}
      \textbf{Axis} & \textbf{C1 -- EV-SpSegNet} & \textbf{C2 -- YOLO+Tracker} & \textbf{C3 -- [TBD]}\\
      \hline
      Accuracy on tiny drones & \textbf{High} (segments curve) & medium (depends on frame stats) & [TBD]\\
      Latency                 & 8\,s window floor & per-frame (\,$\sim$33\,ms) & [TBD]\\
      Hardware                & GPU $+$ \texttt{spconv} & GPU / edge & [TBD]\\
      Failure mode            & static drones (no events) & sparse / slow drones & [TBD]\\
    \end{tabular}
  \end{darkbox}
  \vspace{0.6em}
  {\small \textbf{Plan:} use C1 as cold-start / recovery on top of C2's per-frame tracker;
   C3 fills the remaining gap.}
\end{frame}

% ---------------------------------------------------------------------
% How would we improve (mirrors template)
% ---------------------------------------------------------------------
\begin{frame}{How Would We Improve?}
  \vspace{0.5em}
  \begin{itemize}
    \setlength\itemsep{0.5em}
    \item \textbf{Shrink the window.} Quantify the trade-off between buffer length and IoU --
          does a 4\,s window already give acceptable detection quality?
    \item \textbf{Marry C1 and C2.} Plug the FRED-trained tracker into the
          \texttt{frame\_detector/} bridge and measure end-to-end latency under our hybrid plan.
    \item \textbf{Reuse the STC loss.} Drop $\mathcal{L}_{\text{stc}}$ on top of the FRED frame-based detector
          -- ``penalise isolated positives, reward coherent ones'' applies in 2D too.
    \item \textbf{Address the failure mode.} Stationary / slow-moving drones produce no events
          $\to$ a complementary RGB or IR channel is probably unavoidable (NeRDD, FRED).
  \end{itemize}
\end{frame}

% ---------------------------------------------------------------------
% Thank you
% ---------------------------------------------------------------------
\begin{frame}[plain]
  \vspace{3cm}
  \hspace{1.2cm}%
  {\fontfamily{ptm}\selectfont\fontsize{56}{68}\selectfont Thank You for Listening!}
\end{frame}

\end{document}
